\section{الفصل~\ref{sec:code}. شفرة مورس}

التقديرات الحدية في الفصل من نظرية المعلومات والحساب؛ وما قيس هو أين ينشأ الضجيج في القناة، وما سرعة عمل الدماغ، وكم تكلف النبضة، أما السؤال عن الشيفرة التي تستعملها القشرة فعلًا فيبقى مفتوحًا.

{\footnotesize
\setlength{\tabcolsep}{3pt}\renewcommand{\arraystretch}{1.15}
\begin{longtable}{>{\raggedright}p{0.5cm}>{\raggedright}p{4.0cm}>{\raggedright}p{5.6cm}>{\raggedright}p{2.3cm}>{\raggedright\arraybackslash}p{1.7cm}}
\toprule
م & القضية & التجربة وما قيس & المصدر & الحالة \\
\midrule\endfirsthead
\toprule
م & القضية & التجربة وما قيس & المصدر & الحالة \\
\midrule\endhead
1 & ترددات عصبونات \nt{GLUT} في القشرة من أجزاء النبضة إلى عشرات النبضات في الثانية، والعصبونات تعمل معظم الوقت قرب الحد الأدنى & تسجيل بماصة ملاصقة للخلية (اختيار العصبون لا يتوقف على فعاليته) في القشرة السمعية لجرذ مستيقظ: في كل لحظة يعطي ترددًا فوق $20$ نبضة في الثانية أقل من $5\,\%$ من العصبونات؛ وفي بعض العصبونات يقل التردد الخلفي عن $0.01$ في الثانية؛ وتوزيع الترددات لوغاريتمي طبيعي & \cite{Hromadka2008} & مُقاس \\
2 & سعة القناة النبضية $R_{\max}\approx r\log_2\frac{e}{r\Delta t}$~\eqref{eq:spikeentropy}، وتكاد كلها تكون في توقيت النبضات (القضية~\ref{prop:capacity}) & إنتروبيا سلسلة ثنائية من خانات طولها $\Delta t$. أرقام الفصل: $8$ بتات للنبضة عند $r=10$ في الثانية و$\Delta t=1\ms$، ونحو $5$ بتات عند $\Delta t=10\ms$، وأقل من $2$ بت في الثانية لعدّاد النبضات & \cite{MacKay1952} & نظرية \\
3 & العصبونات الحسية تنقل من بت واحد إلى بضعة بتات للنبضة، وفي أفضل الحالات حتى نصف الحد & استعادة المنبه من نبضات عصبونات حسية معروفة المنبه؛ خلاصة القياسات في الكتاب & \cite{Rieke1997} & مُقاس \\
4 & مولّد النبضات دقيق؛ والتوقيت الدقيق تحدده التغيرات السريعة في المدخل & عصبونات قشرة الجرذ في الشرائح: مع تيار متكرر ذي تذبذبات سريعة تتكرر النبضات بدقة أفضل من $1\ms$، ومع تيار ثابت تتبعثر اللحظات & \cite{Mainen1995} & مُقاس \\
5 & المشبك غير موثوق: أكثر من نصف النبضات لا يصل إلى المستقبِل بسبب عشوائية القذف & التحفيز الأدنى لمداخل العصبونات الهرمية في الحُصين (شرائح): أكثر من نصف النبضات لا يستدعي استجابة. استُبعدت تقلبات عتبة التحفيز، وأعطال التوصيل عند تفرعات المحور، وانخفاض حرارة التجربة؛ فيبقى القذف الاحتمالي & \cite{AllenStevens1994} & مُقاس \\
6 & تدفق النبضات في القشرة الحية يبدو عشوائيًا: الفترات مبعثرة كما في تدفق بواسون، وتباين عدد النبضات قريب من متوسطه؛ وجزء من ”الضجيج“ رسالة عن حركات الحيوان & تسجيلات في المنطقتين البصريتين \foreignlanguage{english}{V1} و\foreignlanguage{english}{MT} لقرد مستيقظ: معامل تغير الفترات نحو $1$، ونسبة تباين عدد النبضات إلى متوسطه كما في عملية عشوائية؛ ونموذج يجمع مداخل صغيرة مستقلة يعطي تشتتًا أصغر بكثير. وفي القشرة البصرية للفأر يُتنبأ بجزء كبير من التشتت من تسجيل الفيديو لحركات الحيوان نفسه & \cite{Softky1993,Stringer2019} & مُقاس \\
7 & ترميز التردد بطيء: الخطأ $1/\sqrt{rT}$~\eqref{eq:rateerror}، والدماغ يتعرف على الصورة في $\sim150\ms$ & الصيغة إحصاء بواسون، استنتاج الفصل. التجربة: قرر أشخاص هل في صورة غير مألوفة معروضة $20\ms$ حيوان؛ تفترق جهود الدماغ لـ”نعم“ و”لا“ بعد نحو $150\ms$. وتقسيم هذا الزمن إلى نحو عشر مراحل من $10\text{--}20\ms$ تقدير الفصل لا قياس & \cite{Thorpe1996} & مُقاس \\
8 & التوسيط على مجموعة يحده الارتباط: لا ينخفض الخطأ بأكثر من $1/\sqrt c$ مرة~\eqref{eq:correlated} (القضية~\ref{prop:pooling}) & أزواج عصبونات من المنطقة \foreignlanguage{english}{MT} لدى القرد مسجلة معًا: معامل ارتباط أعداد النبضات $0.12$ في المتوسط، ويبين التحليل النظري أن حتى هذا الارتباط يحد ربح المجموعة. النظرية: لا يضع الحد إلا الجزء من الارتباط الذي يشبه الإشارة. ونموذج الموقف المقابل: تردد مجموعة من $50\text{--}100$ عصبون يُقرأ في فترة واحدة بين نبضتين، $10\text{--}50\ms$، ولا يُستعمل توقيت النبضات المفردة في القشرة & \cite{Zohary1994,MorenoBote2014,ShadlenNewsome1998} & مُقاس \\
9 & يمكن كتابة العدد في زمن النبضة الأولى~\eqref{eq:latency}، أو في ترتيب نبضات مجموعة، أو في طور الإيقاع المحلي & شبكية السمندل: البنية المكانية لصورة معروضة لحظة مكتوبة في التأخيرات النسبية للنبضات الأولى للعصبونات المخرجة، شبه مستقلة عن التباين. خلايا المكان في حُصين الجرذ: عند عبور موضعها ”الخاص“ تنزاح النبضات أبكر فأبكر في دورة إيقاع ثيتا ($7\text{--}12\,\text{Hz}$)، بإزاحة من $100^\circ$ إلى $355^\circ$. إلى أي حد تستعمل القشرة توقيت النبضات: سؤال مفتوح (السطر~8) & \cite{Gollisch2008,OKeefe1993} & مُقاس \\
10 & أي شيفرة تُقرأ يقرره الثابت الزمني للمستقبِل~\eqref{eq:reader} (القضية~\ref{prop:reader})؛ وفي القشرة النشطة العصبونات أقرب إلى كواشف التزامن & الصيغة~\eqref{eq:reader} استنتاج الفصل. مراجعة لتسجيلات في الكائن الحي: في القشرة النشطة موصلية الغشاء أعلى عدة مرات منها في الراحة، والثابت الزمني أقصر تبعًا لذلك. أما أن القشرة تبدّل الشيفرة بهذه الطريقة تحديدًا فلم يُبيَّن مباشرة & \cite{Destexhe2003} & فرضية \\
11 & المشبك الناضب ينقل تغيّر التردد، أي $\Delta r/r$ في جوهره~\eqref{eq:depression} (الشكل~\ref{fig:depression}) & تسجيلات مزدوجة لعصبونات هرمية في القشرة: سرعة النضوب يحددها احتمال القذف؛ مع النضوب البطيء تعكس الاستجابة التردد، ومع السريع التزامنات. نموذج مبني على النضوب المقيس في قشرة الجرذ: التغيرات النسبية المتساوية في التردد على مداخل سريعة وبطيئة تعطي استجابة متساوية، أي أن المشبك ينقل $\Delta r/r$. والأرقام $1.7\to2.2$ و$6.7$ و$90\ms$ حساب الفصل & \cite{Tsodyks1997,Abbott1997depression} & مُقاس \\
12 & الرشقة ”شَرطة“: احتمال ضياعها $(1-p)^k$~\eqref{eq:burst} صغير، والتيسير يقلله أكثر & مراجعة: في المشابك غير الموثوقة تضيع النبضات المفردة كثيرًا، أما الرشقات فتنتقل بموثوقية بفضل التيسير؛ والرشقة تستدعي تغيرات طويلة الأمد في المشابك. أما أن الدماغ يقرأ الرشقة رمزًا مستقلًا ففرضية & \cite{Lisman1997,AllenStevens1994} & فرضية \\
13 & عصبونات \nt{GABA} السريعة تحدد الإيقاع و”علامات الترقيم“؛ وصنف آخر يثبط المثبِّطات ويفتح البوابة (القضية~\ref{prop:punctuation}) & مراجعة لخصائص أصناف عصبونات \nt{GABA} في القشرة. علم البصريات الوراثي: الإثارة الإيقاعية لعصبونات \nt{GABA} السريعة تستدعي إيقاع غاما، والاستجابة للمنبه تتوقف على طور الدورة. في القشرة البصرية للفأر تثبط عصبونات \foreignlanguage{english}{PV} بعضها بعضًا، وعصبونات \foreignlanguage{english}{SST} كل الأصناف الأخرى، وعصبونات \foreignlanguage{english}{VIP} عصبونات \foreignlanguage{english}{SST} في الغالب. إثارة عصبونات \foreignlanguage{english}{VIP} لدى فأر مستيقظ ترفع التثبيط عن عصبونات \nt{GLUT}؛ وتشغّلها إشارة التعزيز. وتقسيم الأدوار قاعدةً عامة فرضية & \cite{Tremblay2016,Cardin2009,Pfeffer2013,Pi2013} & فرضية \\
14 & تحد الطاقة التردد المتوسط بقيمة من رتبة نبضة واحدة في الثانية~\eqref{eq:energy} & الميزانية الطاقية للمادة الرمادية لدى القوارض من بيانات تشريحية وفيزيولوجية: النبضات والتيارات المشبكية $47\,\%$ و$34\,\%$ من طاقة الإشارات؛ ولا يمكن أن ينشط معًا أكثر من $\sim15\,\%$ من العصبونات. ولقشرة الإنسان: لا يمكن أن ينشط بقوة معًا أكثر من $\sim1\,\%$ من العصبونات. والرقم $\sim10^9$ جزيء ATP للنبضة والقدرة $\sim10\,\text{W}$ للإشارات تقدير الفصل المبني على هذه الحسابات & \cite{Attwell2001,Lennie2003} & نظرية \\
15 & الشيفرة متناثرة ومضبوطة على أكبر عدد من البتات لكل جول؛ والقشرة تقايض الدقة بالطاقة (القضية~\ref{prop:economy}) & فرضية الشيفرة المقتصدة. السند: لدى الفئران عند تقييد الغذاء تهبط في القشرة البصرية موصلية مستقبلات \foreignlanguage{english}{AMPA} والإنفاق المشبكي من ATP ($-29\,\%$)، ويبقى تردد النبضات، ويتسع الضبط على الاتجاه بـ$32\,\%$ & \cite{Barlow1961,Padamsey2022} & فرضية \\
\bottomrule
\end{longtable}}
